% ---
% id: protocol-compiler/rem:source-segment-rules
% kind: rem
% title: "Grounds of the source-segment rules"
% summary: >-
%   The licensed-only rule and the source-mode rule are the two halves of the
%   reductions module's admission rule for a non-empty source segment. Both
%   refuse where the chain is consumed; only the first refuses a case an
%   artefact of this repository can present.
% prominence: standard
% status: accepted
% uses: [def:chaincond]
% mentions: [def:gate-outcome-table, req:gate-outcomes, req:chain-refusals, transmitter/def:source-segment, reductions/def:chain, reductions/lem:precomposition, reductions/rem:ordering, reductions/thm:two-sided, reductions/lem:source-transfer, transmitter/ass:src, transmitter/decoy/cor:emitted-decoy]
% assumes: []
% refs: []
% provenance: newly-derived
% lean:
% ---

Definition~\ref{def:chaincond} states two rules for a non-empty source
segment, the licensed-only rule and the rule on the declared source mode, and
Requirement~\ref{req:chain-refusals} names their refusals. This remark gives
their grounds and compares what each of them refuses.

\begin{remark}\label{rem:source-segment-rules}
\emph{Where the two rules come from.} The licensed-only rule is the consuming
side of the admission rule of \xref{reductions/def:chain}, whose reason is
\xref{reductions/lem:precomposition} and its counterexample
(\xref{reductions/rem:ordering}); this module states neither and decides the
type of the entry. The rule on the declared source mode is the other half of
that admission rule, which admits a non-empty source segment beside a
constraint list only where every entry is licensed \emph{and only in
\tref{trusted-source mode}}: the entry half is the licensed-only rule and the
mode half is this one, and both are decided in
Definition~\ref{def:chaincond} on the presented data, the entry half by
reading the declared tag of each entry and the mode half by reading one
declared slot.

\emph{The licensed-only rule refuses a case the \tref{basic format} admits.}
A non-empty \tref{source segment} is declarable there --- the source-segment
paragraph of Definition~\ref{def:chaincond} admits $L_{\mathrm s}\ge1$ ---
and, \emph{were this rule absent}, an entry declared general beside one
would, where no other condition or gate refused it, reach the gates and be
emitted by row 5 of the gate outcome table of
Definition~\ref{def:gate-outcome-table} as a
declaration that is not \tref{key-length-eligible}
(Requirement~\ref{req:gate-outcomes}). Under the rule the case is refused where the
chain is consumed, before any declaration exists, so row 5 is never reached
for it. The two outcomes are not the same, and the second is strictly the more
conservative. \emph{The case is reachable}, and the rule is then simply a
refusal, which is the direction a consumer may always take.

\emph{Why the source-mode rule refuses rather than weakening the outcome}:
nothing licenses a source map in \tref{untrusted-source mode}. The statements
that make a non-empty segment sound --- \xref{reductions/thm:two-sided} and
\xref{reductions/lem:source-transfer} --- are stated in trusted-source mode
alone, because a source map replaces \emph{Alice's preparation} and only a
trusted specification has one: an untrusted specification carries no
certificate and no marginal, its Alice-side system is an output of the attack,
and A-SRC is asserted in trusted-source mode alone
(\xref{transmitter/ass:src}). So a declaration of that shape would carry a
replaced ensemble with no statement transferring anything to it, which is the
one thing a consumer may not emit. \emph{No artefact of this repository can
present the pairing, and the rule is stated so that it is decided rather than
left to provenance}: a non-empty segment is accepted only inside an artefact
of the module that admits its links, by (A3) of
\xref{transmitter/def:source-segment} as Definition~\ref{def:chaincond}
decides it, and the one submodule that admits a kind
declares its data in trusted-source mode
(\xref{transmitter/decoy/cor:emitted-decoy}\,(i)). Unlike the licensed-only
rule, which refuses a case the basic format declares and which, were the
licensed-only rule absent and nothing else refused it, row 5 would emit as not
key-length-eligible, the source-mode rule
refuses nothing that an artefact of this repository can present; what it
removes is the possibility that a later kind emitted in another mode would
pass every condition stated in Definition~\ref{def:chaincond}.
\end{remark}
